\documentclass[11pt]{article}
\usepackage[letterpaper,margin=1in]{geometry}
\usepackage[T1]{fontenc}
\usepackage{amsmath,amssymb,amsthm,lmodern,microtype}
\usepackage[hidelinks]{hyperref}

\newtheorem{theorem}{Theorem}
\newtheorem{lemma}[theorem]{Lemma}
\newtheorem{proposition}[theorem]{Proposition}

\title{Tuza's Conjecture for Triangle Packing Number at Most Six}
\author{}
\date{}

\begin{document}
\maketitle

\begin{abstract}
Let $\nu(G)$ denote the maximum number of edge-disjoint triangles in a finite simple graph $G$, and let $\tau(G)$ denote the minimum number of edges meeting every triangle of $G$. We prove that every finite simple graph $G$ with $\nu(G)\leq 6$ satisfies Tuza's conjecture $\tau(G)\leq 2\nu(G)$. The proof combines Kandasamy's theorem for graphs of maximum degree at most ten, local link reductions, the Tutte--Berge formula, and the Gallai--Edmonds decomposition.
\end{abstract}

\begin{theorem}\label{thm:main}
Every finite simple graph $G$ with $\nu(G)\leq 6$ satisfies $\tau(G)\leq 2\nu(G)$.
\end{theorem}

The proof uses Kandasamy's theorem for graphs of maximum degree at most ten~\cite{Kandasamy}, the Tutte--Berge formula and Gallai--Edmonds decomposition~\cite{Schrijver}, and local link reductions.

\section{Links in a minimal counterexample}

Suppose Theorem~\ref{thm:main} fails, and choose a counterexample $G$ minimizing $|E(G)|$ and, subject to that, minimizing $|V(G)|$. Every proper subgraph of $G$ has triangle packing number at most six and hence satisfies $\tau\leq 2\nu$.

\begin{lemma}\label{lem:basic}
The graph $G$ has no isolated vertices, every edge of $G$ lies in at least two triangles, and every link $L_v=G[N(v)]$ satisfies $\delta(L_v)\geq 2$. Moreover, for any $v\in V(G)$, any $W\subseteq N(v)$, any nonempty matching $M$ inside $L_v[W]$, and any $Q\subseteq N(v)$ meeting all edges of $L_v$ incident with $W$ outside $M$, we have $|Q|>|M|$.
\end{lemma}

\begin{proof}
An isolated vertex or an edge in no triangle can be deleted without changing $\nu$ or $\tau$. If an edge lies in a single triangle $T$, deleting $E(T)$ reduces $\nu$ by at least one and $\tau$ by at most two (adding the other two edges of $T$ covers every lost triangle), contradicting minimality. Thus every edge lies in at least two triangles, and $\delta(L_v)\geq 2$.

For the scoped link reduction, suppose $|Q|\leq |M|$. Delete the spokes from $v$ to $W$ and the edges of $M$ to form $H$. The $|M|$ triangles $\{v\}\cup e$ ($e\in M$) extend any packing of $H$, while $M$ together with the spokes from $v$ to $Q$ covers every triangle of $G$ missing from $H$ at cost at most $2|M|$, contradicting minimality.
\end{proof}

\begin{lemma}\label{lem:small-links}
We have $\delta(G)\geq 5$, and every link of order five or six is a complete graph minus at most a matching. Consequently, every link has at least seven edges and retains a matching of size two after deleting any two edges.
\end{lemma}

\begin{proof}
A graph of minimum degree at least two on at least four vertices has a matching of size two, because pairwise-intersecting edges form a star or a triangle. A three-vertex link is a triangle, which reduces by Lemma~\ref{lem:basic} using one edge as $M$ and its opposite vertex as $Q$. A four-vertex link reduces using a two-edge matching as $M$ and the endpoints of one matching edge as $Q$. Hence $\delta(G)\geq 5$.

In a five-vertex link $L_v$, if $z$ has two nonneighbors $p,q$, let $Q$ consist of its other two vertices. If $pq\notin E(L_v)$, any two-edge matching gives a reduction with $Q$; if $pq\in E(L_v)$, take $M$ to consist of $pq$ and an edge from $z$ to $Q$. Thus the missing edges form a matching.

In a six-vertex link $L_v$, suppose first that there is no perfect matching. Tutte's theorem gives a separator $S$ whose odd-component count exceeds $|S|$ by at least two, by parity, so $|S|\leq 2$. For $|S|=0$, minimum degree two forces two triangle components, which reduce. For $|S|=1$, minimum degree two prevents isolates, so three odd components cannot fit in five vertices. For $|S|=2$, all four remaining vertices are isolated in $L_v-S$ and adjacent to both vertices of $S$; taking $Q=S$ and a two-edge matching between $S$ and the isolates gives a reduction.

Now suppose $L_v$ has a perfect matching, and some vertex $z$ has two nonneighbors $p,q$. If $d_{L_v}(z)=2$, let $P$ be its three nonneighbors. If $P$ has a nonedge, its ends and $z$ form an independent triple, whose complement serves as $Q$ with any perfect matching. If $P$ is a clique, every perfect matching matches $z$ outside $P$ and therefore uses some edge $p'q'$ of $P$, so $V(L_v)\setminus\{z,p',q'\}$ serves as $Q$. If instead $\delta(L_v)\geq 3$, then $d_{L_v}(z)=3$. Taking $Q=V(L_v)\setminus\{z,p,q\}$ reduces whenever $pq$ is absent or belongs to some perfect matching. Otherwise $L_v-\{p,q\}$ has four vertices, a universal vertex $z$, and no perfect matching, so its other three vertices are independent and their complement serves as $Q$ with any perfect matching of $L_v$. Hence any missing edges form a matching.

Therefore links of orders five and six have at least eight and twelve edges, respectively, and links of order at least seven have at least seven edges by minimum degree two. After deleting any two edges from $L_v$, at least six, ten, or five edges remain for orders five, six, or at least seven, respectively. For orders five and six, this exceeds the maximum edge count of a star or triangle. For order at least seven, a triangle has too few edges, and a star would leave at least six vertices away from its center needing at least six restored incidences to recover minimum degree two, whereas two deleted edges supply at most four.
\end{proof}

\begin{lemma}\label{lem:vc-after-matching}
For every vertex $v\in V(G)$ and every matching $M$ in $L_v$, the graph $L_v-M$ has vertex cover number at least three, and in particular contains a matching of size two.
\end{lemma}

\begin{proof}
Suppose $Q\subseteq V(L_v)$ covers $L_v-M$ with $|Q|\leq 2$. If $|M|\geq 2$, then $M$ and $Q$ give a whole-link reduction (Lemma~\ref{lem:basic}). If $Q$ already covers all of $L_v$, minimum degree two forces $|Q|=2$ and every vertex outside $Q$ adjacent to both vertices of $Q$; since $|V(L_v)|\geq 5$, there is a two-edge matching between $Q$ and $V(L_v)\setminus Q$, giving a reduction with $Q$. The remaining case is a single edge $e\in M$ with both endpoints outside $Q$. Since $|V(L_v)|\geq 5$, choose a vertex $w\in V(L_v)\setminus(Q\cup e)$. By minimum degree two, $w$ has an edge to $Q$, and pairing that edge with $e$ produces a two-edge matching whose deletion leaves all edges covered by $Q$, again a reduction. Finally, any graph with no matching of size two is a star or triangle (plus isolates) and has vertex cover number at most two.
\end{proof}

\section{Matching deficiency at a maximum-degree vertex}

By Kandasamy's theorem~\cite{Kandasamy}, $\Delta(G)\geq 11$. Fix $a\in V(G)$ of maximum degree $d=\Delta(G)$, set $H=G[N(a)]$, and let $X=V(G)\setminus(\{a\}\cup N(a))$. We have $\delta(H)\geq 2$, and $\delta(H)\geq 4$ when $X=\varnothing$. Write $\mu$ for graph matching number.

\begin{lemma}\label{lem:tb}
If a graph $R$ on $d$ vertices satisfies $\mu(R)\leq m$, then there exists $S\subseteq V(R)$ of size $s=|S|\leq m$ such that $R-S$ has $q\geq s+d-2m$ odd components. If $i$ of these odd components are isolated vertices, then
\[
d-s\geq i+3(q-i), \qquad i\geq 2s+d-3m.
\]
Each odd component of $R-S$ has order at least the smallest positive odd integer at least $\delta(R)-s+1$.
\end{lemma}

\begin{proof}
This follows directly from the Tutte--Berge formula~\cite{Schrijver}, noting that $q\leq d-s$ forces $s\leq m$, and $d-s\geq i+3(q-i)\geq 3(s+d-2m)-2i$ yields $i\geq 2s+d-3m$.
\end{proof}

\begin{lemma}\label{lem:isolates}
For any $S\subseteq V(H)$, let $i$ be the number of isolated vertices in $H-S$. If $i\in\{3,4,5\}$, then the isolated vertices supply at least $i+2$ edge-disjoint triangles in $G$, all of whose edges are disjoint from all triangles through $a$ supported on edges of the nonisolated part of $H-S$. In particular, $i\leq 4$, and if $i\in\{3,4\}$ then the nonisolated part of $H-S$ has matching number at most $4-i$.
\end{lemma}

\begin{proof}
The isolated vertices of $H-S$ are pairwise nonadjacent and have all their neighbors in $\{a\}\cup S\cup X$. At the first isolate choose a two-edge matching in its link, and at the second choose a two-edge matching avoiding the two used base edges (Lemma~\ref{lem:small-links}). At each further isolate (up to the fifth), choose one unused link edge; this is possible because each link has at least seven edges and at most six base edges have been used. This yields five, six, or seven edge-disjoint triangles for $i=3,4,5$, respectively. Every edge of these triangles is either a spoke from an isolate to $\{a\}\cup S\cup X$ or a base edge inside $\{a\}\cup S\cup X$. A matching in the nonisolated part of $H-S$ uses vertices disjoint from the isolates and from $\{a\}\cup S\cup X$, so its triangles through $a$ are edge-disjoint from the isolate triangles. Since $\nu(G)\leq 6$, we obtain $i\leq 4$ and the stated bound on the matching number of the nonisolated part.
\end{proof}

\section{The apex must be universal}

\begin{proposition}\label{prop:universal}
The maximum-degree vertex $a$ is universal: $X=\varnothing$.
\end{proposition}

\begin{proof}
Suppose $X\neq\varnothing$. If $\mu(H)\geq 5$, fix a five-edge matching $M$ in $H$, giving five triangles through $a$. If $X$ is independent, take any $x\in X$; then $L_x\subseteq H$, and by Lemma~\ref{lem:vc-after-matching} the graph $L_x-(M\cap E(L_x))$ has a matching of size two, adding two triangles through $x$. If $X$ contains an edge $xy$, choose a triangle $xyz$ containing $xy$; it uses no edge of $H$ and avoids $a$. By Lemma~\ref{lem:vc-after-matching}, $\{y,z\}\cap V(L_x)$ cannot cover $L_x-(M\cap E(L_x))$, so there is an edge $uv\in E(L_x)\setminus M$ avoiding $y$ and $z$. Then $xuv$ and $xyz$ add two triangles disjoint from the five apex triangles. Both cases give $\nu(G)\geq 7$, a contradiction.

Hence $\mu(H)\leq 4$. Apply Lemma~\ref{lem:tb} to $H$ with $m=4$, giving $s\leq 4$, $q\geq s+d-8$, and $i\geq 2s+d-12$.
\begin{itemize}
\item For $s\geq 3$, we have $i\geq 5$, contradicting Lemma~\ref{lem:isolates}.
\item For $s=2$ and $d\geq 13$, again $i\geq 5$. For $s=2$ and $d\in\{11,12\}$, we have $i\geq d-8\geq 3$. If $i\geq 5$, Lemma~\ref{lem:isolates} gives a contradiction. If $i=4$, at least $d-6\geq 5$ nonisolated vertices remain in $H-S$, supplying at least one edge and contradicting Lemma~\ref{lem:isolates}. If $i=3$, then $d=11$ and $q\geq 5$, leaving at least two nontrivial odd components on the remaining six vertices; one edge from each gives a two-edge matching, contradicting Lemma~\ref{lem:isolates}.
\item For $s=1$, minimum degree two rules out isolates, so $3(d-7)\leq d-1$, impossible for $d\geq 11$.
\item For $s=0$, every component has order at least three, so $3(d-8)\leq d$ forces $d\in\{11,12\}$. Then $H$ has at least two three-vertex components (triangles) and a third component containing an edge.
\end{itemize}

For each three-vertex component $uvw$ of $H$, associate to each of its three edges the list of common neighbors of its endpoints in $\{a\}\cup X$. All three lists contain $a$. If two lists equal $\{a\}$, say for $uv$ and $uw$, then $\{a,v,w\}$ is a triangle component of $L_u$, contradicting Lemma~\ref{lem:basic}. If one list equals $\{a\}$ and the union of all three lists has size two, the other two lists contain a common $x\in X$, which is then adjacent to all three vertices and lies in the first list as well. If all three lists equal $\{a,x\}$, then $W=\{a,v,w,x\}$, $M=\{av,wx\}$, and $Q=\{w,x\}$ give a scoped link reduction at $u$ (only $x$ can have link edges leaving $W$). These cases exhaust all failures of Hall's condition. Distinct representatives yield three edge-disjoint triangles for each of the two three-vertex components, and an edge of the third component extends through $a$ to a seventh edge-disjoint triangle, a contradiction.
\end{proof}

\section{A universal apex of degree at least twelve}

Since $X=\varnothing$, we have $V(G)=\{a\}\cup V(H)$ and $\delta(H)\geq 4$. The graph $H$ is connected: otherwise, the cones over its components partition all triangles of $G$ into edge-disjoint proper subgraphs, and summing $\tau\leq 2\nu$ over these cones yields $\tau(G)\leq 2\nu(G)$.

\begin{lemma}\label{lem:odd-cycle}
If a triangle-free graph $F$ on $n$ vertices satisfies $5\delta(F)>2n$, then $F$ is bipartite.
\end{lemma}

\begin{proof}
Suppose $F$ is not bipartite, and let $C$ be a shortest odd cycle, of length $\ell\geq 5$. Then $C$ is chordless, and every vertex $w\notin V(C)$ has at most two neighbors on $C$ (three or more neighbors would be pairwise nonadjacent on $C$, leaving an odd gap of length at most $\ell-4$ that forms a shorter odd cycle with $w$). Counting edges incident with $C$ gives $5\delta(F)\leq\ell\,\delta(F)\leq 2\ell+2(n-\ell)=2n$, a contradiction.
\end{proof}

\begin{proposition}\label{prop:deg12plus}
We have $d=11$.
\end{proposition}

\begin{proof}
Suppose $d\geq 12$. Clearly $\mu(H)\leq 6$. If $\mu(H)\leq 5$, Lemma~\ref{lem:tb} gives $s\leq 5$, $q\geq s+d-10$, and $i\geq 2s+d-15$. For $s\geq 4$, we get $i\geq 5$, contradicting Lemma~\ref{lem:isolates}. For $s=1,2,3$, the minimum odd orders $5,3,3$ multiplied by $s+d-10$ exceed $d-s$ for all $d\geq 12$. For $s=0$, we get $q\geq d-10\geq 2$, contradicting connectivity of $H$.

Hence $\mu(H)=6$. Fix a six-edge matching $M$ in $H$; then $\nu(G)=6$ and $H-M$ is triangle-free. No vertex $v\in V(H)$ has degree four or five in $H$: if $d_H(v)=4$ then $N_H[v]$ induces $K_5$ minus at most a matching, and if $d_H(v)=5$ then $N_H(v)$ induces $K_5$ minus at most a matching, by Lemma~\ref{lem:small-links}. On that five-vertex set, deleting $M$ leaves a graph whose complement is the union of two disjoint matchings and hence bipartite; an independent triple in the complement gives a triangle in $H-M$, a contradiction. Therefore $\delta(H)\geq 6$ and $\delta(H-M)\geq 5$.

If $d=12$, then $5\delta(H-M)\geq 25>24=2d$, so $H-M$ is bipartite by Lemma~\ref{lem:odd-cycle}. The edges of $M$ together with the spokes from $a$ to a smaller bipartition side of $H-M$ cover all triangles of $G$ with at most $6+6=12=2\nu(G)$ edges, a contradiction.

Suppose instead $d\geq 13$. Apply Lemma~\ref{lem:tb} to $H$ with $m=6$ and $\delta(H)\geq 6$, giving $s\leq 6$, $q\geq s+d-12$, and $i\geq 2s+d-18$. For $s=6$, we get $i\geq d-6\geq 7$, contradicting Lemma~\ref{lem:isolates}. For $s=1,2,3,4,5$, the minimum odd orders are $7,5,5,3,3$, respectively; their products with $s+d-12$ exceed $d-s$ at $d=13$ ($14>12$, $15>11$, $20>10$, $15>9$, $18>8$), and each excess increases with $d$. Thus $s=0$, and connectivity of $H$ forces $d-12\leq q\leq 1$, so $d=13$.

If $H$ contains a triangle $T$, let $R=H-E(T)$. Then $\delta(R)\geq 4$, and only the three vertices of $T$ can have degree below six in $R$. If $\mu(R)\leq 5$, Lemma~\ref{lem:tb} gives $s\leq 5$, $q\geq s+3$, and $i\geq 2s-2$. For $s=0,1,2,3$, there are no isolates and the minimum odd orders $5,5,3,3$ require at least $15,20,15,18$ vertices, each exceeding $13-s$. For $s=4,5$, we need at least six or eight isolates, respectively, whereas at most the three vertices of $T$ have degree at most five in $R$. Thus $\mu(R)\geq 6$, which together with $T$ gives seven edge-disjoint triangles in $G$, a contradiction.

Therefore $H$ is triangle-free. Since $5\delta(H)\geq 30>26=2d$, Lemma~\ref{lem:odd-cycle} implies that $H$ is bipartite. The spokes from $a$ to a smaller bipartition side cover all triangles of $G$ using at most six edges, contradicting $\tau(G)>2\nu(G)=12$.
\end{proof}

\section{A universal apex of degree eleven}

\begin{proof}[Proof of Theorem~\ref{thm:main}]
By Propositions~\ref{prop:universal} and~\ref{prop:deg12plus}, $a$ is universal and $H$ is a connected graph on $d=11$ vertices with $\delta(H)\geq 4$.

First we show $\mu(H)=5$. If $\mu(H)\leq 4$, Lemma~\ref{lem:tb} gives $s\leq 4$, $q\geq s+3$, and $i\geq 2s-1$. Connectivity rules out $s=0$. At $s=1$, four odd components of order at least five cannot fit in ten vertices. At $s=2,3$, the forced isolates violate $\delta(H)\geq 4$. At $s=4$, all seven remaining vertices are isolated, contradicting Lemma~\ref{lem:isolates}. Thus $\mu(H)=5$.

Call a vertex of $H$ \emph{exposable} if some five-edge matching of $H$ misses it. We claim every exposable vertex $v$ satisfies $d_H(v)\geq 6$. Let $P$ be a five-edge matching of $H$ missing $v$. If $d_H(v)=4$, then $N_H(v)$ induces $K_4$ minus at most a matching (Lemma~\ref{lem:small-links}) and therefore contains at least two edge-disjoint perfect matchings. Every edge of one meets both edges of the other, so $P$ avoids at least one of them, giving two triangles through $v$ in $H$ disjoint from the five apex triangles of $P$. If $d_H(v)=5$, then $N_H(v)$ induces $K_5$ minus at most a matching; deleting $P$ leaves minimum degree at least two and hence a two-edge matching, again giving two triangles through $v$ avoiding $P$. Both cases give $\nu(G)\geq 7$, a contradiction.

Let $D\subseteq V(H)$ be the set of exposable vertices and $A=N_H(D)\setminus D$, with $s=|A|$. By the Gallai--Edmonds decomposition~\cite{Schrijver}, every component of $H[D]$ is factor-critical (hence odd), and since $H$ has matching deficiency $11-2\mu(H)=1$, there are $s+1$ components of $H[D]$ and $s\leq 5$. Every vertex of $D$ has degree at least six in $H$, with all neighbors in its component or in $A$. For $s=1,2,3,4,5$, the minimum odd component orders $7,5,5,3,3$ multiplied by $s+1$ require at least $14,15,20,15,18$ vertices, each exceeding $11-s$. Thus $A=\varnothing$. Since $D\neq\varnothing$ and $H$ is connected, we conclude $D=V(H)$ and $\delta(H)\geq 6$.

Hence $|E(H)|\geq 33$. If $\nu(H)\leq 1$, minimality of $G$ gives $\tau(H)\leq 2$, and deleting a cover of $H$ leaves a triangle-free graph on eleven vertices with at least $31$ edges, contradicting Mantel's bound $\lfloor 11^2/4\rfloor=30$. Choose two edge-disjoint triangles $T_1,T_2$ in $H$, and set $R=H-E(T_1)-E(T_2)$. Every vertex outside $T_1\cup T_2$ retains degree at least six in $R$, every vertex in exactly one of $T_1,T_2$ has degree at least four, and at most one vertex (shared by $T_1$ and $T_2$) can have degree below four, with degree at least two.

If $\mu(R)\leq 4$, Lemma~\ref{lem:tb} gives $s\leq 4$ and $q\geq s+3$. For $s=0,1$, there are no isolates, and no order-three component can exist because it would require three vertices of degree at most three in $R$; thus every odd component has order at least five, requiring at least $15$ or $20$ vertices, exceeding $11-s$. For $s=2,3$, at most one vertex can be isolated, requiring at least $1+3(q-1)\geq 3s+7>11-s$ vertices. For $s=4$, all seven remaining vertices must be isolated, whereas at most the six vertices of $T_1\cup T_2$ have degree below six in $R$. Therefore $\mu(R)=5$. Combining the five apex triangles from a five-edge matching of $R$ with $T_1$ and $T_2$ yields seven edge-disjoint triangles in $G$, contradicting $\nu(G)\leq 6$.
\end{proof}

\begin{thebibliography}{9}
\bibitem{Kandasamy}
A.~Kandasamy, \emph{Tuza's conjecture for graphs of maximum degree at most 10},
\href{https://arxiv.org/abs/2610.15269}{arXiv:2610.15269}, 2026.

\bibitem{Schrijver}
A.~Schrijver, \emph{Combinatorial Optimization: Polyhedra and Efficiency},
\href{https://www.lamsade.dauphine.fr/~cornaz/Enseignement/M2_MODO/DATA/A3.PDF}{Section~24}, Springer, 2003.
\end{thebibliography}

\end{document}
